Báo cáo trình bày đặc tả kỹ thuật cho \textbf{SmartTraffic What-If (STWI)}, một MVP 13 tuần hỗ trợ người vận hành đánh giá kịch bản giao thông giả định. Hệ thống cung cấp bằng chứng và khuyến nghị; không tự điều khiển đèn tín hiệu hay thiết bị hiện trường.

\bigskip
\textbf{Phạm vi MVP} gồm mạng chức năng 20 node và tối đa 20 luồng camera ghi sẵn/RTSP. Khả năng tiếp nhận quy mô 1.000 camera được kiểm thử bằng producer dữ liệu tổng hợp đã aggregate, không phải 1.000 luồng video inference đồng thời. Video thô được xử lý tại biên và không được lưu mặc định.

\bigskip
Kiến trúc gồm bốn tầng:
\begin{itemize}[leftmargin=1.5cm, itemsep=4pt]
  \item \textbf{Tầng 1 --- Dữ liệu:} YOLOv8 + ByteTrack, MQTT và tensor có node axis $\mathbf{X}\in\mathbb{R}^{B\times12\times N\times16}$, missing mask cùng shape và adjacency $\mathbf{A}\in\mathbb{R}^{N\times N}$.
  \item \textbf{Tầng 2 --- Dự báo và mô phỏng:} GCN--LSTM sinh baseline 30 phút $\mathbf{Y}\in\mathbb{R}^{B\times6\times N\times2}$; surrogate ensemble học từ các lần chạy SUMO offline để đánh giá \textit{IncidentVector}.
  \item \textbf{Tầng 3 --- Tri thức:} Qdrant lưu luật/SOP theo điều khoản và hiệu lực. LLM sinh \textit{SimulationQuery} có kiểu; server dựng SQL tham số hóa trên TimescaleDB bằng read-only role.
  \item \textbf{Tầng 4 --- Điều phối:} LangGraph, Celery, Redis và API job bất đồng bộ. Counterfactual Safety Loop lấy cảm hứng từ CF-VLA nhưng áp dụng chính sách fail-closed và luôn yêu cầu operator phê duyệt.
\end{itemize}

\bigskip
Các mục tiêu nghiệm thu là surrogate $P_{99}<500$ms trên máy 8 CPU cores, 32GB RAM và GPU 12--16GB; E2E $P_{95}\leq30$s và hard deadline/$P_{99}\leq180$s. Khi OOD, uncertainty cao, thiếu citation hoặc vòng an toàn không hội tụ, hệ thống trả \texttt{needs\_review} và không phát hành hành động có thể thi hành.

\bigskip
\noindent\textbf{Từ khóa:} Điều phối giao thông thông minh, GCN--LSTM, Neural Surrogate, SUMO, RAG, Typed Query, Counterfactual Safety Loop, What-if Analysis.
